
\subsection{Chemical Interpretability of Learned Tokens}

To understand what chemical patterns H-Net learns, we analyzed the top 100 most frequent tokens from our best-performing model (PI1M, 22 epochs).

\textbf{Token Categories.} We automatically classified tokens using a hybrid approach combining character pattern heuristics, RDKit validation, and SMARTS functional group matching. Figure~\ref{fig:interpretability}(a,b) shows the distribution: aliphatic patterns (28\%), aromatic rings (20\%), functional groups (14\%), and syntactic elements (11\%).

\textbf{Atom Boundary Respect.} We analyzed whether token boundaries align with atom boundaries in the SMILES string. 69.8\% of tokens fully respect atom boundaries (both start and end), while only 4.1\% split within an atom symbol (e.g., splitting ``Cl'' as ``C'' and ``l''). Figure~\ref{fig:interpretability}(c) compares boundary respect across models.

\textbf{Functional Group Alignment.} We compared H-Net tokens against common functional groups. Simple groups like hydroxyl (-OH) and ethyl (-CC-) are captured as single tokens in $>$99\% of cases. Complex groups like carboxyl (-COOH) and amide (-CONH-) show lower alignment (25-30\%), as H-Net often splits these into multiple semantically meaningful pieces (e.g., ``C(=O)'' and ``N'').

\textbf{Comparison with SmilesPE.} Unlike SmilesPE's chemically-derived vocabulary that achieves higher functional group alignment, H-Net discovers patterns bottom-up from data. While SmilesPE captures benzene rings in $>$99\% of cases as single tokens, H-Net captures only $\sim$58\%---but H-Net learns dataset-specific patterns (like polymer attachment markers ``*)'') not present in SmilesPE's fixed vocabulary.
